\section{SAM3 Agent：VLM + SAM3 的协作架构}

\subsection{架构设计：大脑与眼睛/手}

SAM3 官方代码库提供了一个 Agent 示例，其核心思想是将 VLM 与 SAM3 分工协作：

\begin{importantbox}{SAM3 Agent 的分工模式}
\begin{itemize}
    \item \textbf{VLM（大脑）}：负责高层语义推理、理解用户的自然语言意图、规划执行步骤
    \item \textbf{SAM3（眼睛/手）}：负责底层视觉操作、生成精确的 Mask 分割
\end{itemize}
VLM 不直接画 Mask，而是指挥 SAM3 去画 Mask。
\end{importantbox}

\subsection{工具定义与 System Prompt}

SAM3 Agent 定义了四个工具（Tools），其中最核心的是 \texttt{segment\_phrase}：

\begin{itemize}
    \item 输入：一个简单的名词短语（如 "child"、"blue vest"）
    \item 输出：全图中所有匹配物体的 Mask 图像（带编号标签）
\end{itemize}

System Prompt 中明确规定了关键约束：
\begin{itemize}
    \item \textbf{禁止使用方位词}：不能输入 "left child"
    \item \textbf{禁止使用数量词}：不能输入 "three cats"
    \item \textbf{强制通用名词}：只能用 "child"、"person" 等通用短语
\end{itemize}

\subsection{Visual Grounding 的完整交互流程}

以"找出最左侧穿蓝色马甲的孩子"为例：

\textbf{第 1 轮——分割阶段}：
\begin{enumerate}
    \item GPT-4o 接收原始图像和用户 Query
    \item 因为 SAM3 不擅长方位推理（在 Tool 描述中已说明），GPT-4o 规划策略：先分割出所有穿蓝色马甲的孩子
    \item 调用 \texttt{segment\_phrase("child in blue vest")}
    \item SAM3 返回带编号 Mask 的结构图（检测到 3 个匹配实例）
\end{enumerate}

\textbf{第 2 轮——推理阶段}：
\begin{enumerate}
    \item GPT-4o 接收带 Mask 标注的结构图（这就是 Visual Prompting！）
    \item 进行视觉验证和空间推理，确定哪个 Mask ID 对应"最左侧"的孩子
    \item 输出最终答案
\end{enumerate}

\begin{warningbox}{VLM 的空间推理仍有缺陷}
在实际测试中，GPT-4o 将 Mask ID 1 判断为最左侧，但正确答案应是 Mask ID 2——模型把左右搞混了。这说明即使是最先进的 VLM，在空间位置推理上仍然存在不稳定性。
\end{warningbox}

\subsection{实现细节}

SAM3 Agent 的官方实现并未使用 LangChain/LangGraph 等现代 Agent 框架，而是采用了较为朴素的方式：

\begin{itemize}
    \item 支持文本和图像的混合输入
    \item Tool 定义直接写在 System Prompt 中
    \item 官方示例基于 LLaMA 3 Vision 8B；讲者改用 GPT-4o（支持 Tool + 多模态输入）
\end{itemize}

\subsection{本章小结}

SAM3 Agent 展示了一种经典的"VLM + 专业视觉模型"协作范式：VLM 负责推理与规划，SAM3 负责精准的视觉操作。通过 Tool Call 机制，不具备分割能力的 VLM 获得了像素级别的分割能力，实现了 Visual Grounding。


